\chapter{Fonctions trigonométriques}\label{ch:g12:trigo}

Les \omterm{def:g11:func:function}{fonctions} $\cos$ et $\sin$ transforment la géométrie du cercle en
analyse. Définies en enroulant la droite réelle autour du cercle
trigonométrique, elles sont l'archétype des \omterm{def:g11:func:function}{fonctions} périodiques, et leurs
\omterm{def:g12:deriv:derivative}{dérivées} en font des solutions de l'\omterm{def:g10:algebra:equation}{équation} des oscillations $y'' = -y$.

\section{Le cercle trigonométrique}

\begin{definition}[Cosinus et sinus]\label{def:g12:trigo:cossin}
Soit $\mathcal{C}$ le cercle de rayon $1$ centré à l'origine d'un \omterm{def:g10:coordgeom:system}{repère
orthonormal}. À chaque réel $t$, associer le point $M(t)$ obtenu en
enroulant une longueur $\abs{t}$ le long de $\mathcal{C}$ depuis le point
$I(1,0)$, dans le sens trigonométrique si $t \geq 0$ et dans le sens
horaire sinon. Alors \emph{$\cos t$ et $\sin t$ sont les coordonnées de
$M(t)$}\index{cosinus}\index{sinus}~:
\[
M(t) = (\cos t, \sin t).
\]
\end{definition}

\begin{omfigure}
\begin{tikzpicture}[scale=2.2]
  \draw[->] (-1.25,0) -- (1.35,0) node[right] {$x$};
  \draw[->] (0,-1.25) -- (0,1.35) node[above] {$y$};
  \draw (0,0) circle (1);
  \coordinate (O) at (0,0);
  \coordinate (M) at (50:1);
  \draw[thick,omDef] (O) -- (M);
  \fill (M) circle (0.6pt) node[above right] {$M(t)$};
  \fill (1,0) circle (0.6pt) node[below right] {$I$};
  \draw[dashed] (M) -- (M |- O) node[below,pos=1.05] {$\cos t$};
  \draw[dashed] (M) -- (M -| O) node[left,pos=1.02] {$\sin t$};
  \draw[->,omThm,thick] (0.35,0) arc (0:50:0.35);
  \node[omThm] at (25:0.5) {$t$};
\end{tikzpicture}
\end{omfigure}

Comme le cercle a pour circonférence $2\pi$, le point $M(t)$ est inchangé
lorsque $t$ augmente de $2\pi$~; et comme $M(t)$ est sur le cercle
unité~:

\begin{proposition}[Identités fondamentales]\label{prop:g12:trigo:identities}
Pour tout $t \in \R$ et $k \in \Z$~:
\[
\cos^2 t + \sin^2 t = 1, \qquad
\cos(t + 2k\pi) = \cos t, \qquad
\sin(t + 2k\pi) = \sin t,
\]
\[
\cos(-t) = \cos t, \qquad \sin(-t) = -\sin t,
\]
\[
\cos(\pi - t) = -\cos t, \quad \sin(\pi - t) = \sin t, \quad
\cos\left(\tfrac{\pi}{2} - t\right) = \sin t, \quad
\sin\left(\tfrac{\pi}{2} - t\right) = \cos t .
\]
\end{proposition}

\begin{proof}
La première identité est l'\omterm{def:g10:algebra:equation}{équation} du cercle unité~; les autres expriment
les symétries du cercle~: $M(-t)$ est le symétrique de $M(t)$ par rapport
à l'axe des $x$, $M(\pi - t)$ par rapport à l'axe des $y$, et
$M(\frac\pi2 - t)$ par rapport à la droite $y = x$.
\end{proof}

Les valeurs à connaître par cœur~:
\begin{center}
\begin{tabular}{c|ccccc}
$t$ & $0$ & $\dfrac{\pi}{6}$ & $\dfrac{\pi}{4}$ & $\dfrac{\pi}{3}$ & $\dfrac{\pi}{2}$\\[6pt]
\hline\rule{0pt}{14pt}%
$\cos t$ & $1$ & $\dfrac{\sqrt3}{2}$ & $\dfrac{\sqrt2}{2}$ & $\dfrac12$ & $0$\\[6pt]
$\sin t$ & $0$ & $\dfrac12$ & $\dfrac{\sqrt2}{2}$ & $\dfrac{\sqrt3}{2}$ & $1$
\end{tabular}
\end{center}

\begin{proposition}[Formules d'addition]\label{prop:g12:trigo:addition}
Pour tous $a, b \in \R$~:
\begin{align*}
\cos(a+b) &= \cos a \cos b - \sin a \sin b, &
\sin(a+b) &= \sin a \cos b + \cos a \sin b,\\
\cos(a-b) &= \cos a \cos b + \sin a \sin b, &
\sin(a-b) &= \sin a \cos b - \cos a \sin b.
\end{align*}
En particulier $\cos 2a = 2\cos^2 a - 1 = 1 - 2\sin^2 a$ et
$\sin 2a = 2 \sin a \cos a$.
\end{proposition}

\begin{proof}
Les \omterm{def:g11:vect:vector}{vecteurs} $\vect{u} = (\cos a, \sin a)$ et $\vect{v} = (\cos b, \sin b)$
sont des \omterm{def:g11:vect:vector}{vecteurs} unitaires faisant un angle $a - b$~; leur \omterm{def:g11:scal:dot}{produit
scalaire} égale à la fois $\cos a\cos b + \sin a \sin b$ (\omterm{def:g10:coordgeom:system}{coordonnées}) et
$\norm{\vect u}\,\norm{\vect v}\cos(a-b) = \cos(a - b)$ (géométrie), ce
qui est la troisième formule. Remplacer $b$ par $-b$ donne la première~;
les formules de \omterm{def:g12:trigo:cossin}{sinus} suivent en utilisant
$\sin x = \cos(\frac\pi2 - x)$. Les formules de duplication sont le cas
$b = a$ combiné avec $\cos^2 + \sin^2 = 1$.
\end{proof}

\section{Étude analytique}

\begin{lemma}[Limite fondamentale]\label{lem:g12:trigo:sinxx}
\[
\lim_{t \to 0} \frac{\sin t}{t} = 1
\qquad\text{et}\qquad
\lim_{t \to 0} \frac{\cos t - 1}{t} = 0 .
\]
\end{lemma}

\begin{proof}
Pour $0 < t < \frac{\pi}{2}$, comparer trois aires dans le cercle unité~:
le triangle $OIM(t)$, le secteur circulaire $OIM(t)$, et le triangle
rectangle de base $OI$ et de hauteur $\tan t$~:
\[
\frac{\sin t}{2} \leq \frac{t}{2} \leq \frac{\tan t}{2}.
\]

\begin{omfigure}
\begin{tikzpicture}[scale=2.2]
  \draw[->] (-0.25,0) -- (1.5,0) node[right] {$x$};
  \draw[->] (0,-0.25) -- (0,1.45) node[above] {$y$};
  \fill[omDef!20] (0,0) -- (1,0) -- (50:1) -- cycle;
  \draw[gray, thin] (-10:1) arc (-10:80:1);
  \draw[omProp, thick] (0,0) -- (1,0) -- (1,1.1918) -- cycle;
  \draw[omDef] (0,0) -- (1,0) -- (50:1) -- cycle;
  \draw[omThm, thick] (1,0) arc (0:50:1);
  \draw[dashed] (50:1) -- (0.643,0);
  \node[left, font=\small] at (0.68,0.35) {$\sin t$};
  \node[omProp, right, font=\small] at (1,0.62) {$\tan t$};
  \node[omThm, font=\small] at (20:0.5) {$t$};
  \fill (50:1) circle (0.6pt);
  \node at (58:1.22) {$M(t)$};
  \fill (1,0) circle (0.6pt) node[below right] {$I$};
  \node[below left] at (0,0) {$O$};
\end{tikzpicture}

{\small Les trois aires emboîtées~: le triangle $OIM$ (bleu, aire
$\frac{\sin t}{2}$), le secteur circulaire (borné par l'arc rouge, aire
$\frac{t}{2}$), et le triangle rectangle de hauteur $\tan t$ (orange,
aire $\frac{\tan t}{2}$).}
\end{omfigure}

La première inégalité donne $\frac{\sin t}{t} \leq 1$~; la seconde donne
$\cos t \leq \frac{\sin t}{t}$. Lorsque $t \to 0^+$, $\cos t \to 1$
(\omterm{def:g12:limcont:continuity}{continuité}), et le théorème des gendarmes donne $\frac{\sin t}{t} \to 1$~;
la parité étend cela à $t \to 0$. Pour la seconde limite,
\[
\frac{\cos t - 1}{t} = \frac{\cos^2 t - 1}{t(\cos t + 1)}
= -\frac{\sin t}{t}\cdot\frac{\sin t}{\cos t + 1}
\xrightarrow[t\to0]{} -1 \cdot \frac{0}{2} = 0. \qedhere
\]
\end{proof}

\begin{theorem}[Dérivées du sinus et du cosinus]\label{thm:g12:trigo:derivatives}
Les \omterm{def:g11:func:function}{fonctions} $\sin$ et $\cos$ sont \omterm{def:g12:deriv:derivative}{dérivables} sur $\R$, avec
\[
(\sin)' = \cos, \qquad (\cos)' = -\sin .
\]
Plus généralement, $(\sin(ax+b))' = a\cos(ax+b)$ et
$(\cos(ax+b))' = -a\sin(ax+b)$.
\end{theorem}

\begin{proof}
Par les formules d'addition,
\[
\frac{\sin(x+h) - \sin x}{h}
= \sin x\,\frac{\cos h - 1}{h} + \cos x\,\frac{\sin h}{h}
\xrightarrow[h\to0]{} \sin x \cdot 0 + \cos x \cdot 1 = \cos x,
\]
en utilisant le \cref{lem:g12:trigo:sinxx}. Le calcul pour $\cos$ est
identique, et la forme générale suit de la règle de la chaîne.
\end{proof}

\begin{proposition}[Variations]\label{prop:g12:trigo:variations}
Sur $\intcc{0}{\pi}$, $\cos$ décroît de $1$ à $-1$. Sur
$\intcc{-\frac\pi2}{\frac\pi2}$, $\sin$ croît de $-1$ à $1$. Les deux
\omterm{def:g11:func:function}{fonctions} sont $2\pi$-périodiques et ont pour \omterm{def:g10:functions:function}{image} $\intcc{-1}{1}$~;
$\cos$ est paire, $\sin$ est \omterm{def:g11:func:parity}{impaire}.
\end{proposition}

\begin{proof}
Sur $\intoo{0}{\pi}$, $(\cos)' = -\sin < 0$ car le point $M(t)$ a une
ordonnée positive là~; de même $(\sin)' = \cos > 0$ sur
$\intoo{-\frac\pi2}{\frac\pi2}$. La périodicité et la parité viennent de
la \cref{prop:g12:trigo:identities}.
\end{proof}

\begin{omfigure}
\begin{tikzpicture}
\begin{axis}[
  width=12cm, height=4.5cm,
  axis lines=middle,
  xmin=-7, xmax=7, ymin=-1.4, ymax=1.4,
  xtick={-6.2832,-3.1416,3.1416,6.2832},
  xticklabels={$-2\pi$,$-\pi$,$\pi$,$2\pi$},
  ytick={-1,1},
  samples=200, domain=-7:7]
  \addplot[omDef,thick] {cos(deg(x))};
  \addplot[omThm,thick] {sin(deg(x))};
\end{axis}
\end{tikzpicture}

{\small Les courbes de $\cos$ (bleu) et $\sin$ (rouge).}
\end{omfigure}

\begin{method}[Résoudre $\cos x = a$ et $\sin x = a$]\label{met:g12:trigo:equations}
Pour $a \in \intcc{-1}{1}$, trouver une solution particulière $\alpha$
(d'après le tableau de valeurs ou une calculatrice). Alors toutes les
solutions sont~:
\[
\cos x = \cos\alpha \iff x = \alpha + 2k\pi \text{ ou } x = -\alpha + 2k\pi,
\quad k \in \Z;
\]
\[
\sin x = \sin\alpha \iff x = \alpha + 2k\pi \text{ ou } x = \pi - \alpha + 2k\pi,
\quad k \in \Z.
\]
Pour les inégalités, localiser les arcs solutions sur le cercle
trigonométrique et lire les \omterm{def:g10:numbers:interval}{intervalles} sur une période.
\end{method}

\begin{example}\label{ex:g12:trigo:equation}
Résoudre $\cos x = \frac12$ dans $\intoc{-\pi}{\pi}$~: une solution
particulière est $\frac\pi3$, donc les solutions sont $x = \frac\pi3$ et
$x = -\frac\pi3$. Dans tout $\R$~: $x = \pm\frac\pi3 + 2k\pi$,
$k \in \Z$.
\end{example}

\section{Exercices}

\begin{exercise}[$\star$]\label{exo:g12:trigo:1}
Calculer $\cos\frac{2\pi}{3}$, $\sin\frac{5\pi}{6}$, $\cos\frac{7\pi}{4}$
et $\sin\left(-\frac{\pi}{3}\right)$ en utilisant les identités de la
\cref{prop:g12:trigo:identities}.
\end{exercise}

\begin{exercise}[$\star$]\label{exo:g12:trigo:2}
Dériver $f(x) = \sin^2 x$, $g(x) = \cos(3x) + x\sin x$ et
$h(x) = \dfrac{\sin x}{2 + \cos x}$.
\end{exercise}

\begin{exercise}[$\star$]\label{exo:g12:trigo:3}
Résoudre dans $\intoc{-\pi}{\pi}$, puis dans $\R$~:
\[
\text{(a) } \sin x = \frac{\sqrt3}{2};
\qquad
\text{(b) } \cos\left(2x\right) = \frac{\sqrt2}{2} .
\]
\end{exercise}

\begin{exercise}[$\star\star$]\label{exo:g12:trigo:4}
En utilisant les formules d'addition, calculer la valeur exacte de
$\cos\frac{\pi}{12}$ et $\sin\frac{\pi}{12}$.
(Indication~: $\frac{\pi}{12} = \frac{\pi}{3} - \frac{\pi}{4}$.)
\end{exercise}

\begin{exercise}[$\star\star$]\label{exo:g12:trigo:5}
Résoudre dans $\intcc{0}{2\pi}$ l'inéquation
$2\sin^2 x - \sin x - 1 \geq 0$.
(Indication~: \omterm{def:g10:algebra:expand}{factoriser} le trinôme en $\sin x$.)
\end{exercise}

\begin{exercise}[$\star\star$]\label{exo:g12:trigo:6}
Calculer les limites
\[
\lim_{x\to0} \frac{\sin 3x}{x}, \qquad
\lim_{x\to0} \frac{1 - \cos x}{x^2}, \qquad
\lim_{x\to+\infty} x \sin\frac{1}{x}.
\]
\end{exercise}

\begin{exercise}[$\star\star$]\label{exo:g12:trigo:7}
Étudier la \omterm{def:g11:func:function}{fonction} $f(x) = x + \cos x$ sur $\intcc{0}{2\pi}$~: variations
et extrema. Montrer que $f$ est \omterm{def:g12:seq:monotonic}{croissante} sur $\R$ bien que $f'$
s'annule en une infinité de points.
\end{exercise}

\begin{exercise}[$\star\star\star$]\label{exo:g12:trigo:8}
Soit $f(x) = A\cos x + B \sin x$ avec $(A,B) \neq (0,0)$.
\begin{enumerate}
  \item Montrer que $f$ peut s'écrire $f(x) = R\cos(x - \varphi)$ avec
        $R = \sqrt{A^2 + B^2}$ et un $\varphi$ convenable.
  \item En déduire le \omterm{def:g10:functions:extrema}{maximum} et le \omterm{def:g10:functions:extrema}{minimum} de $f$, et résoudre
        $\cos x + \sin x = 1$ dans $\intoc{-\pi}{\pi}$.
\end{enumerate}
\end{exercise}

\begin{exercise}[$\star\star\star$]\label{exo:g12:trigo:9}
Montrer que $\sin$ et $\cos$ satisfont tous deux l'\omterm{def:g10:algebra:equation}{équation} différentielle
$y'' = -y$. Réciproquement, soit $y$ une solution de $y'' = -y$ avec
$y(0) = a$ et $y'(0) = b$~; montrer que $y = a\cos + b\sin$.
(Indication~: considérer $g = (y - a\cos - b\sin)$ et «~l'énergie~»
$E = g^2 + (g')^2$.)
\end{exercise}

\section{Problème~: le secret du pendule}

\begin{problem}[{Devoir du week-end --- une seule limite géométrique
alimente les dérivées du sinus et du cosinus, et de là toute
l'oscillation~: ressorts, pendules, battements, et le mètre qui a
failli être}]
\label{pb:g12:trigo:1}
Toute horloge qui ait jamais battu, toute corde qui ait jamais
sonné, obéit aux mêmes mathématiques~: une force de rappel
proportionnelle au déplacement, donc l'\omterm{def:g10:algebra:equation}{équation} $y'' = -\omega^2 y$,
donc des \omterm{def:g12:trigo:cossin}{cosinus}. À la base de tout cela se tient une limite d'air
innocent, $\frac{\sin x}{x} \to 1$ --- promise dans le volume
précédent, utilisée par le \cref{thm:g12:trigo:derivatives}, et
démontrée ici par un dessin. Ce problème s'élève de cette limite
jusqu'au balancement d'un pendule d'un mètre, et explique pourquoi
le mètre n'est, à un cheveu près, pas défini par lui.

\medskip
\noindent\textbf{Partie I --- La limite fondamentale.}
\begin{enumerate}
  \item Échauffement numérique~: calculer $\frac{\sin x}{x}$ pour
        $x = 0.1$ et $x = 0.01$ (en \omterm{def:g11:trigo:radian}{radians}~!).
  \item L'encadrement, sur le cercle trigonométrique avec
        $0 < x < \frac\pi2$~: exprimer par de la géométrie
        élémentaire les aires (i) du triangle de sommets $O$,
        $A(1,0)$ et le point du cercle $M(\cos x, \sin x)$~; (ii) du
        secteur circulaire $OAM$~; (iii) du triangle rectangle
        $OAT$, où $T(1, \tan x)$.
  \item De l'encadrement des aires, déduire
        $\sin x < x < \tan x$, puis
        $\cos x < \frac{\sin x}{x} < 1$, et conclure par le théorème
        des gendarmes que $\lim_{x \to 0} \frac{\sin x}{x} = 1$ (la
        parité règle le cas $x < 0$).
  \item En déduire $\lim_{x \to 0} \frac{1 - \cos x}{x^2} =
        \frac12$ (multiplier par la quantité conjuguée).
  \item Expliquer pourquoi cette limite est le moteur du
        \cref{thm:g12:trigo:derivatives} (que vaut $\sin'(0)$~?), et
        pourquoi elle acquitte aussi une vieille dette~: en quel
        sens le \emph{\omterm{def:g11:trigo:radian}{radian}} est-il la seule unité d'angle pour
        laquelle $\sin' = \cos$ tient sans constante parasite~?
\end{enumerate}

\medskip
\noindent\textbf{Partie II --- Le mouvement harmonique.}
\begin{enumerate}[resume]
  \item Vérifier, par dérivation des \omterm{def:g11:func:function}{fonctions} composées, que
        $y(t) = R\cos(\omega t - \varphi)$ satisfait
        $y'' = -\omega^2 y$.
  \item Une masse accrochée à un ressort obéit à $y'' = -4y$, avec
        $y(0) = 3$ et $y'(0) = 8$. En admettant (d'après
        l'\cref{exo:g12:trigo:9}) que toutes les solutions sont de
        la forme $A\cos 2t + B\sin 2t$, déterminer $A$ et $B$.
  \item Réécrire la solution sous la forme $R\cos(2t - \varphi)$
        (\cref{exo:g12:trigo:8})~: donner l'amplitude $R$, la
        période et la phase $\varphi$ (trois décimales).
  \item À quel instant la masse passe-t-elle pour la première fois
        par la position d'équilibre $y = 0$~?
  \item Le pendule~: une masse au bout d'un fil de longueur $L$
        obéit \emph{exactement} à $\theta'' = -\frac gL
        \sin\theta$ --- \omterm{def:g10:algebra:equation}{équation} insoluble par les \omterm{def:g11:func:function}{fonctions}
        élémentaires. Pour de petites oscillations, remplacer
        $\sin\theta$ par $\theta$ (la partie I le justifie~!) et en
        déduire la période $T = 2\pi\sqrt{\frac Lg}$. Calculer $T$
        pour $L = 1$ m (avec $g = 9.81$)~: c'est la célèbre
        découverte de Galilée --- de quoi la formule ne dépend-elle
        \emph{pas}~?
  \item Le pendule battant la seconde~: quelle longueur donne
        exactement $T = 2$ s~? L'Académie des sciences, en 1791,
        hésita entre cette longueur et la dix-millionième partie du
        quart de méridien pour définir le \emph{mètre} --- de
        combien de millimètres les deux candidats
        diffèrent-ils~?
\end{enumerate}

\medskip
\noindent\textbf{Partie III --- Les formules d'addition au
travail.}
\begin{enumerate}[resume]
  \item Montrer avec les formules d'addition
        (\cref{prop:g12:trigo:addition}) que
        $\cos\left(x + \frac\pi2\right) = -\sin x$ et
        $\sin\left(x + \frac\pi2\right) = \cos x$~: dériver, sur le
        cercle, c'est tourner d'un quart de tour --- vérifier la
        cohérence avec le \cref{thm:g12:trigo:derivatives}.
  \item Établir l'identité de transformation de produit en somme
        $\cos a \cos b = \frac12\left(\cos(a - b) +
        \cos(a + b)\right)$ et la linéarisation
        $\cos^2 x = \frac{1 + \cos 2x}{2}$ (le chapitre sur
        l'intégration vous en saura gré).
  \item Battements~: établir
        $\sin p + \sin q = 2 \sin\frac{p + q}{2}
        \cos\frac{p - q}{2}$, et l'appliquer à deux diapasons à
        $440$ et $444$ Hz~: quel son l'oreille entend-elle, et
        combien de pulsations par seconde perçoit-elle~? (C'est
        ainsi que s'accordent les orchestres.)
  \item Établir l'identité de triplication
        $\cos 3x = 4\cos^3 x - 3\cos x$ (écrire $3x = 2x + x$). En
        posant $x = 20^\circ$, quelle \omterm{def:g10:algebra:equation}{équation} du troisième degré
        $\cos 20^\circ$ doit-il vérifier~? (Que cette \omterm{def:g10:algebra:equation}{équation} ne se
        résolve pas à l'aide de seules racines carrées est
        précisément la raison pour laquelle la trisection de l'angle
        a résisté deux millénaires à la règle et au compas ---
        l'histoire complète est racontée dans les volumes
        universitaires.)
  \item Un courant alternatif s'écrit
        $I(t) = 3\cos(100\pi t) + 3\sqrt3 \sin(100\pi t)$. Le
        réécrire sous la forme $R\cos(100\pi t - \varphi)$~:
        intensité de crête et phase~?
\end{enumerate}

\medskip
\noindent\textbf{Partie IV --- Le monde oscillant.}
\begin{enumerate}[resume]
  \item Périodes~: quelle est la (plus petite) période de
        $\sin^2 t$ (utiliser la question 13)~? Et celle de
        $\sin(2t) + \sin(3t)$~?
  \item Un pendule réel s'amortit~:
        $y(t) = \eu^{-t/10}\cos(2\pi t)$. Décrire le mouvement,
        calculer la valeur de l'enveloppe en $t = 10$ (une vieille
        connaissance apparaît, \cref{pb:g12:exp:1}), et nommer le
        chapitre où l'on résout les \omterm{def:g10:algebra:equation}{équations} ayant de telles
        solutions.
  \item Pourquoi les \omterm{def:g12:trigo:cossin}{sinus} et les \omterm{def:g12:trigo:cossin}{cosinus} règnent-ils sur toute
        vibration~? Donner en deux phrases le dictionnaire
        physique-mathématiques (force de rappel proportionnelle au
        déplacement $\rightarrow$ quelle \omterm{def:g10:algebra:equation}{équation} $\rightarrow$
        quelles solutions), et dire ce que
        l'\cref{exo:g12:trigo:9} apporte à cette affirmation.
  \item Pour finir~: l'arc du chapitre en quatre pas --- une limite
        géométrique, deux \omterm{def:g12:deriv:derivative}{dérivées}, une \omterm{def:g10:algebra:equation}{équation} différentielle,
        toute l'oscillation. Et l'astérisque honnête~: pour de
        grandes oscillations, la vraie période du pendule fait
        intervenir une intégrale elliptique, calculée à la vitesse
        de l'éclair par\,\dots\ la \omterm{def:g11:stat:mean}{moyenne} arithmético-géométrique
        de Gauss, le dessert du \cref{pb:g12:seq:1}. Les fils de la
        série se nouent d'eux-mêmes.
\end{enumerate}
\end{problem}
